\documentclass[a4paper]{article}
\usepackage[pdftex]{graphicx}
\usepackage[utf8]{inputenc}
\usepackage{enumerate}
\usepackage{siunitx}
\sisetup{locale=DE}
\usepackage{href-ul}
\hypersetup{
	colorlinks=true,
	linkcolor=blue,
	urlcolor=blue}
\usepackage{icomma}
\usepackage{amssymb}
\usepackage{geometry}
\geometry{a4paper, top=15mm, left=15mm, right=15mm, bottom=15mm,
	headsep=10mm, footskip=12mm}

\begin{document}
	%Title
	{\bf School assignment from physics }\\
	\begin{enumerate}[1.]
		%Task 1
		{\bf \item In an experiment, the change in length of a rod was measured as a function of temperature.} The following series of measurements resulted:
		
		\begin{tabular}{cccccccc}
			$\vartheta$ to $\SI{}{\celsius}$ & 0 & 10 & 15 & 20 & 28 & 35 & 40\\
			$l$ in $\SI{}{\milli\meter}$ & 1000.0 & 1000.3 & 1000.6 & 1000.7 & 1000.9 & 1001.3 & 1001.5
		\end{tabular}
		
		\begin{enumerate}[a)]
			\item Now evaluate the experiment graphically by drawing a $\Delta l$ -- $\Delta \vartheta$ -- diagram.
			\item Evaluate the series of measurements numerically and calculate the length expansion number
			$\alpha$.
		\end{enumerate}
		
		%Exercise 2
		{\bf \item A steel tank with volume $V=60~l$ is completely filled with gasoline at a temperature of $\SI{15}{\celsius}$.} How much gasoline overflows when the temperature rises to $ \SI{40}{\celsius}$ increased?\\
		$[ \gamma_{Gasoline}= \SI{0.00106}{\kelvin^{-1}} \quad
		\gamma_{Steel}= \SI{0.000036}{\kelvin^{-1}}$
		
		%Task 3
		{\bf \item A cold liquid thermometer is quickly immersed in hot water.}
		One observes that the liquid column first goes down the scale, then turns around and rises until the thermometer finally shows the actual temperature value. How can this initial downward softening be explained?
		
		%Task 4
		\begin{minipage}{0.6\textwidth}
			{\bf \item In a closed pipe system with the cross section shown in the sketch there is water at the temperature $\SI{7}{\celsius}$.} Now leg A is wrapped with ice so that the water there initially reaches $ \SI{4}{\celsius}$ cools down.
			\begin{enumerate}[a)]
				\item Now explain what can be observed.
				\item What happens if the water continues to cool down to $\SI{1}{\celsius}$?
			\end{enumerate}
		\end{minipage}
		\hfill
		\begin{minipage}{0.35\textwidth}
			\includegraphics[width=4 cm]{zirk106.pdf}
		\end{minipage}
		
		%Task 5
		{\bf \item task}
		\begin{enumerate}[a)]
			\item What is Gay-Lussac's law about the behavior of volume and temperature of ideal gases and under what conditions is it valid?
			(The law must be stated in words and as a formula)
			\item What does Boyle-Mariotte's law say about the pressure and volume of an ideal gas? Derive the general gas equation from these two laws.
		\end{enumerate}
		
		{\bf \item A weather balloon rises from the ground to a height of 10,000 m.} The following conditions apply:
		
		\begin{tabular}{ccc}
			& floor & height \\
			Temperature & $\SI{18}{\celsius} $ & $\SI{-50}{\celsius}$ \\
			Pressure & $\SI{1000}{\hecto\pascal} $ & $\SI{264}{\hecto\pascal}$ \\
			Volume & $\SI{1,2}{\meter^3}$ & \\
		\end{tabular}
		
		Calculate the volume of the balloon at altitude.
		
		{\bf \item Energy is supplied to a closed gas by means of radiation. }
		\begin{enumerate}[a)]
			\item Initially the volume of the gas remains constant. What can be said about the change in the kinetic energy $E_{kin}$ and the potential energy $E_{pot}$ of the particles of the gas?
			\item As the gas progresses, it expands while maintaining its temperature (e.g. a piston in a cylinder is pushed outwards by the gas). What now applies to the change in the $E_{kin}$ and the $E_{pot}$ of the particles?
		\end{enumerate}
		
	\end{enumerate}
	
	\fbox{
		\begin{minipage}{0.5\textwidth}
			For the solution please \href{https://www.okuyakl.de/phys/p9thdL106/le106.pdf}{click here} or scan the QR code.\\
			Additional worksheets are available at
			
			\href{http://www.okuyakl.com}{www.okuyakl.com}
		\end{minipage}
		\hfill
		\begin{minipage}{0.4\textwidth}
			\includegraphics[width=1.5 cm]{../../viecher/zwe03}
			\includegraphics[width=3 cm]{qre106}
			\includegraphics[width=2 cm]{../../viecher/afanticon1}
			
	\end{minipage}}
	
\end{document}